\section{The lifting theorem}\label{sec:lifting}

The construction proceeds through finite-dimensional central division
algebras. At each step we impose an algebraic relation on one prescribed
element while preserving the freedom to enlarge the algebra at the next step.
Formal power series in two noncommuting variables record that freedom.

For a $k$-algebra $A$, write
\[
 A\langle\!\langle z_1,z_2\rangle\!\rangle
\]
for the ring of formal series in words in $z_1,z_2$, with coefficients
in $A$ commuting with the variables. Multiplication concatenates words
and multiplies their coefficients in their original order. The ring is
complete for the filtration by word length. For a tuple
$X=(X^{(1)},\ldots,X^{(u)})$ of such series, let $X_0\in A^u$ be its
constant tuple and let $X_1,X_2\in A^u$ be its coefficients of the two
individual letters. Thus the subscripts $1,2$ denote letters, not degrees.
Set
\begin{equation}\label{eq:inner-tangent}
 J_A(X_0)=\{([T,X_0^{(1)}],\ldots,[T,X_0^{(u)}]):T\in A\},
 \qquad [T,a]=Ta-aT.
\end{equation}
The classes of $X_1,X_2$ modulo this subspace describe changes of the tuple
that cannot be obtained by infinitesimal conjugation.

\begin{theorem}[Algebraic lifting]\label{thm:lifting}
Let $k$ be a countable field of characteristic zero containing every root
of unity, and let $A$ be a finite-dimensional central division algebra
over $k$. Suppose
$X\in A\langle\!\langle z_1,z_2\rangle\!\rangle^u$ satisfies:
\begin{enumerate}[label=\textup{(\roman*)}]
\item $X_0$ generates $A$ as a $k$-algebra;
\item $X_1,X_2$ have linearly independent classes in $A^u/J_A(X_0)$.
\end{enumerate}
For every $w\in A\langle\!\langle z_1,z_2\rangle\!\rangle$, there exist
a countable field extension $k'/k$, a central division algebra $A'/k'$ of
strictly larger degree, and series
\[
 X'\in A'\langle\!\langle s_1,s_2\rangle\!\rangle^u,
 \qquad w'\in A'\langle\!\langle s_1,s_2\rangle\!\rangle
\]
with the following properties:
\begin{enumerate}[label=\textup{(\alph*)}]
\item Every noncommutative polynomial identity over $k$ satisfied by the
tuple $(X,w)$ is satisfied by $(X',w')$ as an identity of formal series.
\item The constant tuple $X'_0$ generates $A'$ over $k'$, and the classes
of $X'_1,X'_2$ in $(A')^u/J_{A'}(X'_0)$ are linearly independent.
\item There is a monic nonzero polynomial $q\in k'[t]$ such that $q(w')=0$.
\end{enumerate}
Here the degree of a central division algebra $B/F$ is $\sqrt{[B:F]}$.
\end{theorem}

Theorem~\ref{thm:lifting} is the central construction. Its application to
the Kurosh problem uses $u=2$ and addresses one polynomial in the two
generators at a time. The fields are allowed to grow; they need not be
algebraic extensions of their predecessors. The final step will therefore
identify the actual center separately and prove infinite dimension over it.

Figure~\ref{fig:response-interfaces} separates the two tasks within a
lifting step. The response calculation makes the prescribed series
algebraic; the final translation restores the hypotheses needed to repeat
the step. The two transverse directions must survive together, not merely
one at a time.

\begin{figure}[tbp]
\centering
\begin{tikzpicture}[
  x=1cm,y=1cm,>=Stealth,
  every node/.style={font=\small},
  interface/.style={draw=black!65,rounded corners=2pt,align=center,
    text width=54mm,minimum height=11mm,inner sep=5pt},
  endpoint/.style={interface,text width=117mm,fill=black!3},
  flow/.style={draw=black!65,->,line width=.8pt}
]
  \node[endpoint] (input) at (0,0)
    {$X_0$ generates $A/k$; two independent classes in $A^u/J_A(X_0)$\\[3pt]
     \textcolor{blue!65!black}{$[X_1]$}\qquad
     \textcolor{orange!80!black}{$[X_2]$}};
  \node[interface] (responses) at (-3.15,-1.75)
    {Loop and paired eigenvalue changes\\
     Section~\ref{sec:responses}};
  \node[interface] (errors) at (3.15,-1.75)
    {First unsolved word errors\\
     Section~\ref{sec:word-transfer}};
  \node[interface] (linkage) at (3.15,-3.3)
    {Functional linkage and replication\\
     Section~\ref{sec:linkage-replication}};
  \node[interface] (relations) at (-3.15,-3.3)
    {Common spectral relations\\
     Section~\ref{sec:global-defect}};
  \node[interface] (cyclic) at (-3.15,-4.85)
    {Cyclic compatibility of word errors\\
     Section~\ref{sec:cyclic-obstructions}};
  \node[interface] (correction) at (3.15,-4.85)
    {Eleven-block correction: $q(w(Z))=0$\\
     Section~\ref{sec:free-correction}};
  \node[endpoint] (output) at (0,-6.7)
    {Division realization and translation
     (Sections~\ref{sec:division-realization}--\ref{sec:generation})\\[3pt]
     $X'_0$ generates $A'/k'$; two independent classes in
     $(A')^u/J_{A'}(X'_0)$\\[3pt]
     \textcolor{blue!65!black}{$[X'_1]$}\qquad
     \textcolor{orange!80!black}{$[X'_2]$}};
  \draw[flow] (input.south) -- ++(0,-.3) -| (responses.north);
  \draw[flow] (responses.east) -- (errors.west);
  \draw[flow] (errors.south) -- (linkage.north);
  \draw[flow] (linkage.west) -- (relations.east);
  \draw[flow] (relations.south) -- (cyclic.north);
  \draw[flow] (cyclic.east) -- (correction.west);
  \coordinate (outputbend) at ($(correction.south)!0.5!(output.north)$);
  \draw[flow] (correction.south) |- (outputbend) -| (output.north);
\end{tikzpicture}

\caption{The spectral case of Theorem~\ref{thm:lifting}.
The two displayed classes at each end are linearly independent in the
indicated quotient. The middle boxes describe the matrix calculation:
allowable eigenvalue changes control the first unsolved word errors,
and the cyclic relations are precisely the obstructions to correction.
Realization and translation return a generating tuple over a division
algebra of strictly larger degree. If $w$ already satisfies the minimal
polynomial of $w(0)$, the middle spectral correction is omitted, but
realization and translation remain necessary. No embedding tower of
the successive division algebras is asserted.}
\label{fig:response-interfaces}
\end{figure}
\clearpage
